%%%%%%%%%%%%%%%%%%%%%%%%%%%%%%%%%%%%%%%%%
% FRI Data Science_report LaTeX Template
% Version 1.0 (28/1/2020)
% 
% Jure Demšar (jure.demsar@fri.uni-lj.si)
%
% Based on MicromouseSymp article template by:
% Mathias Legrand (legrand.mathias@gmail.com) 
% With extensive modifications by:
% Antonio Valente (antonio.luis.valente@gmail.com)
%
% License:
% CC BY-NC-SA 3.0 (http://creativecommons.org/licenses/by-nc-sa/3.0/)
%
%%%%%%%%%%%%%%%%%%%%%%%%%%%%%%%%%%%%%%%%%


%----------------------------------------------------------------------------------------
%	PACKAGES AND OTHER DOCUMENT CONFIGURATIONS
%----------------------------------------------------------------------------------------
\documentclass[fleqn,moreauthors,10pt]{ds_report}
\usepackage[english]{babel}

\graphicspath{{fig/}}




%----------------------------------------------------------------------------------------
%	ARTICLE INFORMATION
%----------------------------------------------------------------------------------------

% Header
\JournalInfo{FRI Natural language processing course 2026}

% Interim or final report
\Archive{Project report} 
%\Archive{Final report} 

% Article title
\PaperTitle{Natural language processing course project: Retrieval-Augmented Generation for the Gaia dataset} 

% Authors (student competitors) and their info
\Authors{Guillem Masdemont Serra, Pietro Sestito, and Plabon Shaha}

% Advisors
\affiliation{\textit{Advisors: Slavko Žitnik}}

% Keywords
\Keywords{}
\newcommand{\keywordname}{Keywords}


%----------------------------------------------------------------------------------------
%	ABSTRACT
%---------------------------------------------------------------------------------------

\Abstract{This project aims at enahancing a larger public to get information from one of the most vaste and complex astronomy datasets, the Gaia one, represnting the Milky Way. In order to do so, we plan to implement a retrieval-augmented generation (RAG) system. Such system allows the users to get information through natural language rather than through the use of the Astronomic data query language (ADQL), a SQL-based language.}
%----------------------------------------------------------------------------------------

\begin{document}

% Makes all text pages the same height
\flushbottom 

% Print the title and abstract box
\maketitle 

% Removes page numbering from the first page
\thispagestyle{empty} 

%----------------------------------------------------------------------------------------
%	ARTICLE CONTENTS
%----------------------------------------------------------------------------------------

\section*{Introduction}
	
The aim of this project is to facilitate the access of a larger general public to the Gaia dataset, a vast and complex relational database on the Milky Way \cite{GaiaMission2016, 2022A&A...667A.148G}.
Due to the high dimension and complexity of the topic and the dataset, the main current way of getting information from it is through the use of the Astronomic data query language (ADQL), a SQL-based language \cite{GAVIP2016}. 
In order to make the data more accessible, we propose to implement a retrieval-augmented generation (RAG).
Such system allows to retrieve relevant information from external data sources, properly process it and finally combine the initial query and the retrieved information to get to a final output, often through a large language model.
One of the main advantages of such a choice is that it allows the users to get information through natural language. More specifically, considering a representative subsample of the dataset, whose whole dimension is around 10 TB, we will implement a RAG system able to answer to several different kinds of questions.
The main areas of interest are the spatial distribution, the galactic depth analysis, the photometric characterization and the kinematics of the solar neighborhood.



%------------------------------------------------

\section*{Background Studies}
Large Language Models (LLMs) store knowledge in their parameters, and have revolutionized the field of Natural Language Processing (NLP). However, they suffer from a number of limitations, including the inability to access up-to-date information, the tendency to produce hallucinated content, and they perform poorly on tasks that require tool-use. To tackle these issues, two new paradigms have gained popularity, namely Retrieval Augmented Generation (RAG) and Tool-augmented LLMs. 
Here we present a brief overview of these two paradigms:

\subsection*{Retrieval Augmented Generation (RAG)}
The RAG system is a combination of three fundamental components:

\begin{itemize}
	\item \textbf{Retrieval} - Responsible for retrieving relevant information from external datasources (e.g. a knowledge base, a search engine, a vector database, etc.) based on the input query. The component can be further divided into multiple subcomponents, including query reformulation, decomposition, re-packing, re-ranking, etc. 
	\item \textbf{Augmentation} - The component is responsible for processing, extracting and summarizing the retrieved information, and preparing it to be fed to the generator.
	\item \textbf{Generator} - Combines the input query and the retrieved information to produce the final output. The generator is typically a large language model, but it can also be a smaller model or even a rule-based system.
\end{itemize}

\subsection*{Tool Calling and Agentic LLMs}
Tool-augmented LLMs are a paradigm in which the LLM is augmented with the ability to call external tools (e.g. a calculator, a search engine, a database, etc.) to perform specific tasks. The LLM can be trained to learn when and how to call the tools, and how to integrate the results from the tools into its output. This paradigm is particularly useful for tasks that require reasoning, multi-step problem solving, or access to up-to-date information. For example, an LLM can query a database of astronomical observations or compute physical quantities instead of approximating them through text.

Agentic LLMs are a special case of tool-augmented LLMs, in which the LLM is trained to learn how to call multiple tools and how to integrate the results from the tools into its output. The agentic LLM can be trained to learn how to call the tools in a specific order, and how to use the results from the tools to perform complex tasks.
%------------------------------------------------

\section*{Proposed project dataset/corpus/data/..}

The Gaia mission, launched by the European Space Agency has the objective of creating the largest and most precise three-dimensional map of our Galaxy, the Milky Way, by surveying approximately 1\% of its 100 billion stars \cite{GaiaMission2016, 2022A&A...667A.148G}. The most recent full dataset (2022), Gaia Data Release 3 (DR3) expanded this catalog to include radial velocities for 33 million stars and spectra for 220 million sources \cite{GaiaDR32023}. By the end of 2026, researchers expect the release of the Gaia Data Release 4 (DR4), which is projected to include more than 2.7 billion stars \cite{GaiaDR4_Cosmos}. This upcoming release will significantly increase the dimensionality of the dataset by providing more than $100$ observations for all sources \cite{GaiaScanningLaw2012}. Consequently, the information storage requirements are expected to surge from the 10 TB of DR3 to more than 500 TB of processed data for DR4 alone \cite{GaiaDR4_Projected_2026, Guy2022} eventually reaching petabyte scale by the end of the decade \cite{Vagg2016}. 


All Gaia data releases are made accessible to the public through the official \href{https://archives.esac.esa.int/gaia}{Gaia Archive} \cite{GaiaDR32023}. The information is structured within a complex relational database, where users retrieve specific data using the SQL-based Astronomical Data Query Language (ADQL) \cite{GAVIP2016}. 

Despite this open-access policy, the sheer volume and technical complexity of the dataset often present significant barriers to entry for non-professional audiences and astronomy enthusiasts \cite{Fabricius2021}. While substantial effort has been directed toward leveraging Gaia for educational outreach, the large volume of the releases is often too difficult for clasroom setting or casual exploration. 

The primary objective of this project is to implement a Retrieval-Augmented Generation (RAG) system to bridge this gap. By utilizing a natural language interface, this system aims to empower the public to perform sophisticated queries, making the Gaia dataset more accessible for both educational enrichment and independent scientific inquiry.

In this project, we propose to utilize the \textbf{Gaia Data Release 3 (DR3) dataset}, which is publicly available via the \href{https://archives.esac.esa.int/gaia}{official Gaia Archive}. For now, we will restrict our study to a representative subsample of the catalog and our primary focus is the construction of an automated pipeline designed to resolve the following research queries, which were purposed by \textbf{Prof. Mercè Romero} from the University of Barcelona, a rellevant institution within the Gaia Data Processing and Analysis Consortium (DPAC). 

\begin{enumerate}
    \item \textbf{Spatial Distribution:} Which stars are located within a specific celestial region? 
    This task involves executing a \textit{cone search} by defining a central coordinate (Right Ascension, Declination) and a search radius. The output will be a curated subset of Gaia sources within the specified solid angle.
    
    \item \textbf{Galactic Depth Analysis:} What is the effective limiting distance of observations toward the Galactic Center versus the North Galactic Pole? This requires a directional cone search integrated with distance estimatestable to account for parallax inversions and systematic uncertainties.
    
    \item \textbf{Photometric Characterization:} What is the color distribution of the stellar population in a given direction? 
    By retrieving the color index from the \texttt{gaia\_source} table, the pipeline can generate photometric histograms to analyze stellar temperatures and evolutionary stages.
    
    \item \textbf{Kinematics of the Solar Neighborhood:} At what velocities do the stars nearest to the Sun rotate?
    This involves selecting sources with high parallax values (indicating proximity). These parameters allow for the computation of total space velocity relative to the local standard of rest. 
\end{enumerate}




%----------------------------------------------------------------------------------------
%	REFERENCE LIST
%----------------------------------------------------------------------------------------
\bibliographystyle{unsrt}
\bibliography{report}


\end{document}